\documentclass[11pt]{scrartcl}

\usepackage[stixtwo]{fontsetup}
\usepackage{amsmath}

\begin{document}
$a_2 = \sqrt{10} > a_1$, és ha $a_n < a_{n + 1}$, akkor
$6a_n - 8 < 6a_{n + 1} - 8$, tehát
\(
a_{n + 1} = \sqrt{6a_n - 8} < \sqrt{6a_{n + 1} - 8} = a_{n + 2}
\),
vagyis a sorozat monoton nő.

Ha a sorozat korlátos, akkor konvergens, és határértéke a sorozat legkisebb
felső korlátja. A lehetséges $A$ határérték kielégíti az
$A = \sqrt{6A - 8}$ egyenletet, vagyis $A = 2$ vagy $A = 4$. Mivel
$a_n \ge 3$, így a határérték csak $4$ lehet.

Belátjuk, hogy $a_n < 4$: $a_1 < 4$, és ha $a_n < 4$, akkor $6a_n - 8 < 16$,
így $a_{n + 1} = \sqrt{6a_n - 8} < 4$.

A sorozat tehát monoton és korlátos, így konvergens, a határértéke $4$.
\end{document}
